\documentclass[14pt,a4paper]{extarticle}
\usepackage[T2A]{fontenc}\usepackage[utf8]{inputenc}\usepackage[russian]{babel}
\usepackage{geometry}\geometry{left=30mm,right=15mm,top=20mm,bottom=20mm}
\setlength{\parindent}{1.25cm}\setlength{\parskip}{0.3em}\sloppy
\begin{document}\begin{titlepage}\begin{center}
САНКТ-ПЕТЕРБУРГСКИЙ НАЦИОНАЛЬНЫЙ\\ ИССЛЕДОВАТЕЛЬСКИЙ УНИВЕРСИТЕТ ИТМО
\vspace{1.5cm}

Дисциплина: Бэк-энд разработка
\vspace{1.5cm}

\textbf{Отчёт}\\[0.5cm]
Лабораторная работа 7\\[0.5cm]
\textbf{Облачное развёртывание веб-приложения}
\vfill
\begin{flushright}Выполнил:\\ Цатинян Артём\\ Группа БР1.2\\[0.8cm] Проверил:\\ Добряков Д. И.\end{flushright}
\vfill Санкт-Петербург\\ 2026 г.\end{center}\end{titlepage}
\section{Задача и состояние}
Спроектировать развёртывание приложения в облаке в соответствии с ДЗ7, описать инфраструктуру как код и автоматическую доставку. Ниже приведён проект развёртывания: облачный кластер и пайплайн в рамках настоящего отчёта не создавались и не проверялись.
\section{Исходное приложение}
Приложение бронирования ресторанов включает сервисы identity, catalog, booking и review, PostgreSQL и RabbitMQ. Локальная конфигурация приложения использует Docker Compose; в ЛР5 работоспособность проверялась на отдельной виртуальной машине через nginx. Эти результаты не являются проверкой облачного развёртывания.
\section{Проект инфраструктуры}
В ДЗ7 предложен Yandex Cloud: управляемый Kubernetes, PostgreSQL, реестр контейнеров и балансировщик. Внешний трафик через балансировщик и Ingress поступает в сервисы Kubernetes, затем в приложения. Доступ к базе ограничивается сетью и правилами безопасности, учётные данные хранятся отдельно от репозитория. RabbitMQ предполагается разместить в кластере с постоянным томом; единственный экземпляр является точкой отказа.
\section{Инфраструктура как код}
Предлагается описать сеть, кластер, группы узлов, базу и необходимые права с помощью Terraform. Переменные окружений и секреты не включаются в репозиторий. Приложения предлагается описать Deployment и Service либо Helm-чартом, включая probes, requests и limits. Перед применением необходимо проверить доступность ресурсов и стоимость.
\section{Сборка и доставка}
Пайплайн должен проверять код, собирать образы с идентификатором коммита, загружать их в реестр и применять версионированную конфигурацию Kubernetes. Затем проверяется готовность подов и HTTP-доступность каждого сервиса. При ошибке выполняется откат к предыдущей версии.
\section{План проверки}
Проверить состояние Terraform, узлов и подов, доступ к API через Ingress, соединения с БД и брокером; выполнить перезапуск пода и повторную проверку. Фактических результатов и снимков облачной инфраструктуры пока нет.
\section{Вывод}
Сформирован план облачного развёртывания; практическая реализация и проверка остаются открытыми.
\end{document}
